\section{4 сентября 1914 г.}

\lp{20260604_193811850}

Швейцария. Цюрих. \\
Emile Medtner \\
Zurich \\
Hotel Bellevue

\lp{20260604_193822658}

Дорогой мой! Сегодня подали вашу телеграмму, что вы получили три письма. Хоть
три, и то слава Богу. Написала я массу, т.к. писала каждый день понемногу, а
отправляла через день. Марго будет у нас послезавтра, тогда поговорю с ней о
Мусагете. Она очеь занята своим лазаретом. У нее 17 человек пока. Еще привезут.
Каждый день которому-нибудь делают операцию, и она всегда присутствует. Наш
\rem{раненный} совсем почти поправился, и его вероятно скоро заменят другим. Я
думала держать одного раненного вплоть до

\lp{20260604_193834816}

\noindent
вашего возвращаения (он у прислуги в комнате, а не в Колиной) и
предупредила, что возьму может быть ненадолго. Теперь, когда вы собираетесь
остаться там всю зиму, я пока что не упраздню наш маленький лазарет. О делах вы
не беспокойтесь. Мусагет до вашего возвращения продержится. У Коли кое-как
начинают собираться уроки. Два новых \rem{ученика}, но только по 15 рублей за
два урока ввиду военного времени. А старые ученики все где-то застряли, кто за
границей, кто в обстоятельствах.

Коля страшно обрадовался, что до вас дошли наконец хоть три письма, и теперь
тоже напишет вам, а то его парализовала мысль, что письма совсем не доходят.

\lp{20260604_193834816}

Впрочем, одно письмо он вам послал. Сейчас уже час ночи, а потому бросаю и
говорю вам «покойной ночи», мой милый.

\section{5 сентября 1914 г.}

\lp{20260604_193834816}

Это лепестки роз, которые увянут сейчас у меня в комнате. Дойдет ли до вас их
запах? Сейчас почти 2 часа ночи. Покойной ночи, мой дорогой!

Струве призвали. Он так и не дождался своей семьи, которая застряла в Берлине.
Теперь жена приедет, а его нет. Не могу писать вам о том что мы сейчас
переживаем здесь.
